
\subsubsection{3.1 Hypothesis Structure}

The research question is formulated as a structured hypothesis with testable predictions:

\textbf{Core Hypothesis (H-StaticFeedback-v1):} Under iterative LLM code repair on HumanEval/MBPP, integrating static analyzer feedback (pylint/mypy) alongside execution feedback improves pass@k compared to execution-only feedback, because static analysis provides mechanistic error explanations.

\textbf{Causal Mechanism:}

\begin{enumerate}
\item Static analyzer produces warnings on LLM-generated code
\item LLM receives mechanistic feedback (why, not just where)
\item Targeted repairs lead to higher pass@k
\end{enumerate}

\textbf{Key Assumptions:}

\begin{itemize}
\item A1: Static analyzers produce meaningful warnings on LLM code
\item A2: LLMs can interpret static warnings into fixes
\item A3: Static and execution feedback catch non-overlapping errors
\item A4: HumanEval contains problems with static-detectable errors
\item A5: Pylint severity filtering (disabling conventions) is effective
\end{itemize}

\subsubsection{3.2 Gate-Based Validation}

A hierarchical gate structure is employed to validate assumptions before full experiments:

\begin{table}[htbp]
\centering
\resizebox{\columnwidth}{!}{%
\begin{tabular}{llll}
\toprule
Gate & Type & Hypothesis & Criterion \\
\midrule
h-e1 & MUST\_WORK & Existence & $\geq 30$\% of problems have actionable warnings \\
h-m1 & MUST\_WORK & Mechanism & LLM interprets warnings into fixes \\
h-m2 & SHOULD\_WORK & Mechanism & Static/execution non-overlapping \\
h-c1 & SHOULD\_WORK & Comparison & Improvement stratified by warnings \\
\bottomrule
\end{tabular}
}
\end{table}

MUST\_WORK gates block downstream experiments if they fail. This design catches methodology issues early before resources are expended on invalid experiments.

\subsubsection{3.3 Experimental Design}

\textbf{Variables:}
\begin{itemize}
\item Independent: Feedback type (execution-only vs. static+execution)
\item Dependent: pass@1, pass@5, pass@10
\item Controlled: Base LLM, iteration budget (5), prompt format, static analyzer (pylint)
\end{itemize}

\textbf{Dataset:} HumanEval (164 problems)

\textbf{Static Analysis Configuration:}
\begin{verbatim}
PYLINT_ARGS = ["--output-format=json", "--disable=C,R"]
PYLINT_TIMEOUT_SEC = 30
ACTIONABLE_TYPES = ("error", "warning")
GATE_THRESHOLD = 0.30
\end{verbatim}

Convention (C) and Refactoring (R) categories are disabled; only error and warning types are considered actionable.

\subsubsection{3.4 Existence Gate Protocol}

Before full pass@k evaluation, the existence assumption is validated: do static analyzers produce actionable warnings on benchmark code?

\textbf{Protocol:}
\begin{enumerate}
\item Load all HumanEval problems (164)
\item Extract code (canonical solutions or LLM generations)
\item Run pylint with configured filters
\item Count problems with at least one actionable warning
\item Gate passes if warning\_rate $\geq$ 30\%
\end{enumerate}

\textbf{Threshold rationale:} Below 30\%, fewer than 1 in 3 problems would receive static feedback, limiting both statistical power and practical utility. This threshold is conservative; lower rates might enable meaningful experiments but would require larger sample sizes.

\subsubsection{3.5 Metrics}

\begin{itemize}
\item \textbf{Warning Rate:} Fraction of problems with $\geq 1$ actionable warning
\item \textbf{Warnings per Problem:} Average actionable warnings across all problems
\item \textbf{Warning Distribution:} Breakdown by pylint warning code
\end{itemize}
